% Options for packages loaded elsewhere
\PassOptionsToPackage{unicode,linktoc=all}{hyperref}
\PassOptionsToPackage{hyphens}{url}
%
\documentclass[
]{article}
\usepackage{amsmath,amssymb}
\usepackage{lmodern}
\usepackage{iftex}
\ifPDFTeX
  \usepackage[T1]{fontenc}
  \usepackage[utf8]{inputenc}
  \usepackage{textcomp} % provide euro and other symbols
\else % if luatex or xetex
  \usepackage{unicode-math}
  \defaultfontfeatures{Scale=MatchLowercase}
  \defaultfontfeatures[\rmfamily]{Ligatures=TeX,Scale=1}
\fi
% Use upquote if available, for straight quotes in verbatim environments
\IfFileExists{upquote.sty}{\usepackage{upquote}}{}
\IfFileExists{microtype.sty}{% use microtype if available
  \usepackage[]{microtype}
  \UseMicrotypeSet[protrusion]{basicmath} % disable protrusion for tt fonts
}{}
\makeatletter
\@ifundefined{KOMAClassName}{% if non-KOMA class
  \IfFileExists{parskip.sty}{%
    \usepackage{parskip}
  }{% else
    \setlength{\parindent}{0pt}
    \setlength{\parskip}{6pt plus 2pt minus 1pt}}
}{% if KOMA class
  \KOMAoptions{parskip=half}}
\makeatother
\usepackage{xcolor}
\usepackage[top=3cm,left=3cm,right=3cm,bottom=3cm]{geometry}
\usepackage{graphicx}
\makeatletter
\def\maxwidth{\ifdim\Gin@nat@width>\linewidth\linewidth\else\Gin@nat@width\fi}
\def\maxheight{\ifdim\Gin@nat@height>\textheight\textheight\else\Gin@nat@height\fi}
\makeatother
% Scale images if necessary, so that they will not overflow the page
% margins by default, and it is still possible to overwrite the defaults
% using explicit options in \includegraphics[width, height, ...]{}
\setkeys{Gin}{width=\maxwidth,height=\maxheight,keepaspectratio}
% Set default figure placement to htbp
\makeatletter
\def\fps@figure{htbp}
\makeatother
\setlength{\emergencystretch}{3em} % prevent overfull lines
\providecommand{\tightlist}{%
  \setlength{\itemsep}{0pt}\setlength{\parskip}{0pt}}
\setcounter{secnumdepth}{5}
\ifLuaTeX
  \usepackage{selnolig}  % disable illegal ligatures
\fi
\usepackage{csquotes}
\IfFileExists{bookmark.sty}{\usepackage{bookmark}}{\usepackage{hyperref}}
\IfFileExists{xurl.sty}{\usepackage{xurl}}{} % add URL line breaks if available
\urlstyle{same} % disable monospaced font for URLs
\hypersetup{
  pdftitle={Thermal2 User Manual},
  hidelinks,
  pdfcreator={LaTeX via pandoc}}

\title{Thermal2 User Manual}
\usepackage{etoolbox}
\makeatletter
\providecommand{\subtitle}[1]{% add subtitle to \maketitle
  \apptocmd{\@title}{\par {\large #1 \par}}{}{}
}
\makeatother
\subtitle{User Manual and Quick Reference}
\author{}
\date{}

\makeatletter
\def\@xobeysp{\mbox{}\space}
\def\verbatim@font{\normalfont\ttfamily\raggedright}
%\def\enquote#1{{``#1''}}
\makeatother

%lp
\setmainfont{TeX Gyre Schola}
\setmathfont{TeX Gyre Schola Math}

\begin{document}
\begin{titlepage}
\begin{center}
\def\version{1.2.0}
\def\qeImage{images/qe.png}


  \includegraphics[width=5cm,height=2cm,keepaspectratio]{\qeImage}
  \hspace{2cm} 
  \includegraphics[width=5cm,height=2cm,keepaspectratio]{images/anharmonic.png}
  
	       \vspace{5.5cm}
  \Huge Thermal2 User Manual (v.\version) 
  \vspace{5.5cm}

  \includegraphics[width=3cm,height=1.5cm,keepaspectratio]{images/sciences-sorbonne.png}
  \hspace{2cm} 
  \includegraphics[width=3cm,height=1.5cm,keepaspectratio]{images/impmc.png}
  \hspace{2cm} 
  \includegraphics[width=3cm,height=1.5cm,keepaspectratio]{images/cnrs.png}

\end{center}

\end{titlepage}

{
\setcounter{tocdepth}{3}
\tableofcontents
}
\hypertarget{foreword}{%
\section{Foreword}\label{foreword}}

The thermal2 suite of codes has been written starting in 2014 by Lorenzo
Paulatto{[}\protect\hyperlink{ref1}{1}{]}. It descends from an initial
set of unreleased codes written by Giorgia
Fugallo{[}\protect\hyperlink{ref1}{1},\protect\hyperlink{ref2}{2}{]} and
Andrea Cepellotti{[}\protect\hyperlink{ref3}{3}{]}, both have also
participated in the development. The Wigner conductivity expression has
been developed by Michele
Simoncelli{[}\protect\hyperlink{ref17}{17},\protect\hyperlink{ref18}{18}{]}.

\hypertarget{copyright}{%
\subsection{Copyright}\label{copyright}}

All the files are provided under the
\href{https://www.gnu.org/licenses/old-licenses/gpl-2.0.en.html}{GPL
license, v2} or newer and, when possible, under the
\href{https://cecill.info/licences/Licence_CeCILL_V2.1-fr.html}{CeCILL
license}. A single file nist\_isotopes\_db.f90 contains public domain
data from the National Institute of Standards and Technology.

\hypertarget{citing}{%
\subsection{Citing}\label{citing}}

We would greatly appreciate if when using the thermal2 suite of codes
you cite the following papers where the underlying theory is described
in detail (see also the \protect\hyperlink{Bibliography}{the
bibliography}):

\begin{itemize}
\tightlist
\item
  All applications: L. Paulatto, F. Mauri, and M. Lazzeri, Phys. Rev.~B
  87, 214303 (2013)
\item
  Exact BTE solution: G. Fugallo, M. Lazzeri, L. Paulatto, and F. Mauri,
  Phys. Rev.~B 88, 045430 (2013)
\item
  Spectral functions: L. Paulatto, I. Errea, M. Calandra, and F. Mauri,
  Phys. Rev.~B 91, 054304 (2015)
\item
  Finite size effects: L. Paulatto, D. Fournier, M. Marangolo, M.
  Eddrief, P. Atkinson, M. Calandra. Phys. Rev.~B 101 (20), 205419
  (2020)
\item
  For Wigner conductivity: Simoncelli, Marzari, Mauri, Nature Physics
  volume 15, 809--813 (2019) and Simoncelli, Marzari, Mauri, Physical
  Review X 12 041011 (2022).
\end{itemize}

\hypertarget{table-of-contents}{%
\section{Table of contents}\label{table-of-contents}}

\hypertarget{quick-reference-manual}{%
\section{Quick Reference Manual}\label{quick-reference-manual}}

\hypertarget{compiling-the-code}{%
\subsection{Compiling the code}\label{compiling-the-code}}

The thermal2 codes come bundled with the D3Q code, used to compute
ab-initio the 3rd order dynamical matrices. We refer you to the D3Q
manual for more detailed instruction on compiling. A separate
stand-alone distribution of the thermal2 codes is being considered and
may be available in the future.

\hypertarget{list-of-codes}{%
\subsection{List of codes}\label{list-of-codes}}

The thermal2 suite of codes contains a number of small specialized codes
that only take command line arguments, and a few larger codes that read
input from a file. This is the list of codes and a brief description.

\hypertarget{main-codes}{%
\subsubsection{Main codes}\label{main-codes}}

Each of these codes will be reviewed in detail in a separate section.

\hypertarget{harmonic-phonons-d3_r2q.x}{%
\paragraph{Harmonic phonons
(d3\_r2q.x)}\label{harmonic-phonons-d3_r2q.x}}

This code can compute the phonon frequency at a q-point or along a path,
it can optionally print out the dynamical matrix. It is analogous to the
matdyn.x code of the Quantum-ESPRESSO (QE) Phonon suite, but it uses the
highly optimized subroutines developed for d3\_lw.x and d3\_tk.x. It
should be used to check the phonon dispersion before doing more serious
calculations. It can also compute some useful harmonic phonon
quantities: 1. group velocities 2. interpolated dynamical matrices at
any q-point 3. mean square displacement of atoms at a give temperature
4. phonon internal energy and zero-point energy

See also \protect\hyperlink{d3_r2qx}{d3\_q2r.x}.

\hypertarget{quasi-harmonic-approximation-d3_qha.x}{%
\paragraph{Quasi harmonic approximation
(d3\_qha.x)}\label{quasi-harmonic-approximation-d3_qha.x}}

An efficient and very easy to use quasi-harmonic approximation
implementation. It includes hydrostatic pressure effect and
equation-of-state fitting. See also
\protect\hyperlink{d3_qhax}{d3\_qha.x input format}

\hypertarget{third-order-linewidth-d3_lw.x.}{%
\paragraph{Third order linewidth
(d3\_lw.x).}\label{third-order-linewidth-d3_lw.x.}}

This code computes phonon anharmonic properties: 1. phonon linewidth
(i.e.~its inverse lifetime, aka HWHM, aka the Imaginary part of the
bubble-diagram self energy), 2. the entire bubble-diagram self energy
(comprising both the third order linewidth and the lineshift) and 3. the
phonon spectral weight for a range of energies. 4. the final state
decomposition of an energy, interpreted as two-phonons scattering

See also \protect\hyperlink{d3_lwx}{d3\_lw.x input format}.

It always includes the intrinsic anharmonic contribution from
phonon-phonon interaction and can optionally include Casimir border
scattering and isotope scattering.

\hypertarget{thermal-conductivity-d3_tk.x}{%
\paragraph{Thermal conductivity
(d3\_tk.x)}\label{thermal-conductivity-d3_tk.x}}

The tk code computes the thermal conductivity. It can use the
Single-Mode Approximation (SMA) or the variational approaches
implemented by Fugallo et.al. {[}2{]} using a robust conjugate gradient
minimization. It can include Casimir and isotope scattering {[}15{]}.
See also \protect\hyperlink{d3_tkx}{d3\_tk.x input format}.

\hypertarget{temperature-dependent-phonons-d3_tdph.x}{%
\paragraph{Temperature-dependent phonons
(d3\_tdph.x)}\label{temperature-dependent-phonons-d3_tdph.x}}

A simple reciprocal space implementation of the temperature-dependent
effective potential (TDEP) method. Still experimental, only second order
force constants are include! See also
\protect\hyperlink{d3_tdphx}{d3\_tdph.x input format}.

\hypertarget{thermal2-utilities}{%
\subsubsection{Thermal2 utilities}\label{thermal2-utilities}}

\hypertarget{d3_q2r.x}{%
\paragraph{d3\_q2r.x}\label{d3_q2r.x}}

This code is analogous to the q2r.x code of QE, and it uses the same
input, but produces a file of Force Constants (FCs) which has already
been re-centered in the reciprocal space Wigner-Seitz cell to make
Fourier Interpolation faster.

In addition to the standard q2r variables, you can specify
\enquote{nfar} which is the distance from the origin, in unit cells,
used to construct the first Brilloouin zone. Setting nfar=2 produces
\enquote{centered} force constants, useful for Fourier interpolation.
Setting nfar=0 produces standard \enquote{periodic} force constants that
can be directly compared with the ones from Quantum ESPRESSO.

\hypertarget{d3_qq2rr.x}{%
\paragraph{d3\_qq2rr.x}\label{d3_qq2rr.x}}

Analogous to q2r.x, but operates on the 3rd order matrices. This codes
takes as command line arguments the dimension of the q-points grid and
optionally the name of the output file. You must feed it feed by
standard input the list of anharmonic dynamical matrices in the XML
format produced by d3q.x. For example, d3q was run with
fild3dyn=\enquote{anh} for a NQX × NQY × NQZ grid, you can compute the
3rd order FCs as:

\begin{verbatim}
ls anh* | qq2rr.x NQX NQY NQZ [-o mat3R] [-f NFAR] [-w]
\end{verbatim}

The result will be written by default to a file called mat3R, but you
can change this with the -o option. You can specify the optional
argument NFAR (default=2) which is the number of neighbors that the code
will scan to build the Wigner-Seitz cell. The default value usually
works, a larger value may be necessary with very anisotropic cells or if
atoms are placed very far away from the origin. You can set it to zero
to avoid any re-centering: The resulting force constant will be
unsuitable for doing any further calculation, due to aliasing in the
Fourier interpolation, but they may be easier to import/convert to
another code format.

The code will automatically select from the list the files that it needs
to fill the grid; all the other files will be discarded. Which means
that you can easily compute the force constants for any sub-lattice of
the one available, just by changing the values of NQX, NQY and NQZ.

After writing the force constant to file, the code will perform two
optional tests (you can skip them pressing CTRL-C). First test: for this
test the initial D3 matrices will be recomputed using the force
constants with both the real and imaginary parts (which should be zero).
Second test: recompute the D3 matrices with only the real part of the
force constants. If any discrepancy is detected it will be printed on
output. Notice that any discrepancy in the first test indicate a very
serious problem with the D3 calculation. On the other hand, some
discrepancy is inevitable in the second test; especially if you your
atoms where not in the theoretical equilibrium positions. Also,
increasing the cutoff and k-points can improve the consistency of the
second test.

If the -w option is specified, when performing the FFT test, if the
re-computed D3 matrix differs significantly from the initial one it will
be printed to a file. The file will start with prefix
\enquote{anh\_cmplx} for the first test and \enquote{anh\_real} for the
second test.

\hypertarget{d3_sparse.x}{%
\paragraph{d3\_sparse.x}\label{d3_sparse.x}}

This code converts a file of third order FCs from dense form to sparse
form; it can optionally discard elements that are smaller than a custom
threshold. It can also measure the speedup gained by using the sparse
FCs instead of the dense ones. Syntax:

\begin{verbatim}
sparse.x [-i mat3R.input] [-o mat3R.output] 
        [-t threshold] [-n num_trials]
\end{verbatim}

Where ma3R.input (default mat3R) is the name of the dense file of Fcs,
mat3R.output will be the output file of sparse FCs (you can use
\enquote{none} to avoid saving them to file); threshold is in
Ry/bohr{[}3{]} (all matrix elements smaller than this will be discarded,
default: zero, do not discard anything) and num\_trials is the number of
random trial q-point triplets to compute by Fourier interpolation. If
num\_trials is provided, the code will print out the elapsed time using
the dense and sparse algorithm, the speedup and the eventual discrepancy
between the two methods (which should be zero if the threshold is zero)

\hypertarget{d3_asr3.x}{%
\paragraph{d3\_asr3.x}\label{d3_asr3.x}}

This code applies the acoustic sum rules (ASR) to the third order FCs.
It can only work on dense Fcs, not on the sparse ones. As the sum is
applied iteratively, it will automatically stop after 10,000 iterations,
or when the residual violation of the ASR is less than 10{[}-12{]} or if
a file named STOP is found in the running directory. Syntax:

\begin{verbatim}
asr3.x [-i mat3R.input] [-o mat3R.output]
       [-t threshold] [-n iter_max]
\end{verbatim}

These options will read the dense FCs from file mat3R.input, apply the
ASR iteratively until threshold is reached (default 10{[}-12{]}), or for
iter\_max, then save it to mat3R.output (default: mat3R.input.asr). If a
file named \enquote{STOP} is found in the working directory, the code
will stop after the next iteration and immediately save the FCs to
mat3R.output

\hypertarget{d3_recenter.x}{%
\paragraph{d3\_recenter.x}\label{d3_recenter.x}}

NOTE: this code is useful for debugging, but it is provided \enquote{as
is}, with no support or guarantee.

\begin{verbatim}
d3_recenter.x NQX NQY NQZ [-n NFAR]
           [-i mat3R.input] [-o mat3R.output] [-w]
\end{verbatim}

Reads force constants from mat3R.input, interpolate them on a grid of
NQX × NQY × NQZ points, recenter them on a Wigner-Seitz cell constructed
up to NFAR unit cells and save the result in mat3R.input.recenter.

Uses the properties of Fourier interpolation to transform the 3rd order
force constants from a grid to another. If the new grid is different
than the initial one, some interpolation will be done, if the grid is
the same, you can use the nfar parameter to recalculate the Wigner-Seitz
cell centering. This code be useful to compare the results from grids of
different sizes, or to put the force constants in a format that is
easier to understand for external codes.

If the -w option is specified, the intermediate D3 matrices will, for
the NQX × NQY × NQZ grid will be written to files called
atmp\_Q1\ldots\_Q2\ldots\_Q\ldots.

\hypertarget{d3_import_shengbte.x}{%
\paragraph{d3\_import\_shengbte.x}\label{d3_import_shengbte.x}}

\begin{verbatim}
d3_import_shengbte.x NQX NQY NQZ [-n NFAR] [-w] [-s mat2R]
           [-i FORCE_CONSTANT_THIRD] [-o mat3R.shengbte] 
\end{verbatim}

Reads the 3-body force constants produced by Mingo \& Carrete code
thirdorder.py6 and import them to the thermal2 format. The size of the
supercell used for the FCs calculation must be specified as NQX × NQY ×
NQZ. In order to prepare the FCs for Fourier interpolation, they are
taken to reciprocal space and then back to real space, and re-centered
including up to NFAR neighbouring cell to ensure locality.

In order to read the system information (cell and position of the atoms)
a file containing the force constants header in the thermal2 format must
be provided with the -s option (default: mat2R). Note that a 2nd order
FCs file, produced with d3\_q2r.x or even with normal q2r.x, for the
same system, is sufficient.

If the -w option is specified, the intermediate D3 matrices, generated
on the NQX × NQY × NQZ grid, will be written to files with names
atmp\_Q1\ldots\_Q2\ldots\_Q3\ldots{} (check the manual of d3q for
details on the file names).

See also \protect\hyperlink{import_phonopypy}{import\_phonopy.py}.

\hypertarget{d3_sqom.x}{%
\paragraph{d3\_sqom.x}\label{d3_sqom.x}}

NOTE: this code is experimental and not widely tested, use at your own
risk.

This codes reads a spectral weight file from d3\_lw.x and computes the
convolution with a Lorentzian function that has an energy-dependent
FWHM. This procedures simulates the broadening introduced by Raman
spectroscopy experiments. This code can also sum and average the
spectral function coming from several different files, to simulate the
uncertainty of the neutron wavevectors. It reads its input from a file
called input.SQOM. Please see teh example input.SQOM in Examples for
details.

\hypertarget{experimental-codes}{%
\subsubsection{Experimental codes:}\label{experimental-codes}}

\hypertarget{d3_db.x}{%
\paragraph{d3\_db.x}\label{d3_db.x}}

Uses the ansatz of ref. \protect\hyperlink{ref16}{16} to apply the
anharmonic correction to the dynamical matrix (instead that on the
phonon mode). Can be used to obtain 3rd-order corrected matrices that
can be interpolated. In takes mostly the same input variables as
d3\_lw.x in \protect\hyperlink{calculation-character-default-lw}{the
\enquote{lw full} case} (but in namelist \&dbinput) and will print out
dynamical matrix files for all the requested q-points.

\hypertarget{tools}{%
\subsubsection{Tools}\label{tools}}

In the tools/ directory, a selection of tools for pre- and
post-processing of data.

\hypertarget{funcoft.sh}{%
\paragraph{funcoft.sh}\label{funcoft.sh}}

This is a short bash script to get a plottable file of the linewidth of
a specific phonon mode as a function of temperature.

It reads a list of linewidth files produced by d3\_lw.x and extracts the
frequency, linewidth (and if possible shifted frequency) from all files
for a specific phonon q-point and band and prints a list ordered by
temperature and smearing. Syntax:

\begin{verbatim}
  funcoft.sh point mode file [file2 [file3...]]
\end{verbatim}

\begin{itemize}
\tightlist
\item
  point: index of the point to collect among the files
\item
  mode: number of the phonon mode 1\ldots3~nat
\end{itemize}

The output will contain 5 columns:

\begin{enumerate}
\def\labelenumi{\arabic{enumi}.}
\tightlist
\item
  temperature
\item
  smearing
\item
  omega
\item
  linewidth (gamma, HWHM)
\item
  shifted frequency (if available in the file)
\end{enumerate}

This script is provided \enquote{as is}, using non-standard names for
the output files can break it.

\hypertarget{recompute-sma.m}{%
\paragraph{recompute-sma.m}\label{recompute-sma.m}}

A simple octave/mathlab script that allow you to recompute the thermal
conductivity in single mode approximation using the output from d3\_tk.x
(using store\_lw=.true.) and d3\_r2q.x (using
calculation=\enquote{extr}). This script allow you to quickly combine
different intrinsic/extrinsic scattering sources without repeating the
entire calculation, to manually change parameter and to extract useful
information, like the per-mod contribution to thermal transport. Please
note that this is not a brainless script: some editing (i.e.~at the very
list the unit cell volume) and understanding of the physics is requires.

\hypertarget{diffd3.x}{%
\paragraph{diffd3.x}\label{diffd3.x}}

Compare two D3 files, write on output the matrix elements and the
maximum difference. Take as arguments either two file names (the first
and second D3 matrix) or two directory names a file name, which will be
opened in both directories.

\hypertarget{d3_sc2c.x}{%
\paragraph{d3\_sc2c.x}\label{d3_sc2c.x}}

Open two D3 files, the first for a unit cell calculation for an
arbitrary triplet, the second for a supercell calculation for a triplet
of kind (0,q,-q). Then it refold the super-cell D3 matrix to the unit
cell and compares the two.

The following scripts are in the tools subdirectory, they can be useful
in specific circumstances. They have little documentation, do not
hesitate to ask for help if you cannot make them work.

\hypertarget{xml2giorgia.x}{%
\paragraph{xml2giorgia.x}\label{xml2giorgia.x}}

Reads a list of D3 matrix files in XML format from standard input and
write them to a single ASCII file called d3.txt

\hypertarget{apply_asr.sh}{%
\paragraph{apply\_asr.sh}\label{apply_asr.sh}}

\begin{verbatim}
apply_asr.sh [-i FILDYN.in] [-o FILDYN.out] [-a ASR_TYPE]
\end{verbatim}

A simple bash script that applies the sum rule to a set of dynamical
matrix files (FILDYN.in\emph{, default: dyn) produced by phonon and
saves them with a different name (FILDYN.out}, default: asr\_dyn).
Useful to apply the sum rule \enquote{crystal} (default for ASR\_TYPE)
which is not supported by the thermal2 codes yet. The final fildyn files
can be used normally with d3\_q2r.x or q2r.x.

\hypertarget{import_phonopy.py}{%
\paragraph{import\_phonopy.py}\label{import_phonopy.py}}

In tools you can also file a python script import\_phonopy\_qe.py and
import\_phonopy\_vasp.py to import the FORCE\_CONSTANTS files of 2-body
force constants produced by phonopy (with QE and VASP respectively).
This script is in a very rudimentary stage, it will produce a file
called \enquote{fc} with the force constants in the standard QE format.
You will have to convert it to the thermal2 format with the fc2mat2R.sh
script, also found in tools. A similar tool to convert 3-body force
constants from phono3py could be obtained generalizing this one, but
lack of documentatio makes it difficult to test.

\hypertarget{fc2mat2r.sh}{%
\paragraph{fc2mat2R.sh}\label{fc2mat2r.sh}}

Convert a force constants file from the standard q2r format to the
optimized format produced by d3\_q2r and used by thermal2. Do not use
this script if you can use the original dynamical matrix files from
phonon, use them instead as it is more accurate.

\hypertarget{apply_asr.sh-1}{%
\paragraph{apply\_asr.sh}\label{apply_asr.sh-1}}

Apply the acoustic sum rule to a set of dynamical matrices using q2r and
matdyn. Thermal2 does not implement the more sofisticate sum-rule
methods, but you can apply them directly to the dynamical matrices using
this little script.

\hypertarget{codes-input-and-output}{%
\section{Codes input and output}\label{codes-input-and-output}}

The three main codes read their configuration from input files. All the
variables included in the namelists can also be specified as command
line arguments with the format

\begin{verbatim}
--keyword value
\end{verbatim}

You may have to put quotes around the value if it contains special
characters or spaces. Note that this mechanism is still a bit
experimental, and there is currently no way to specify the data blocks
(QPOINTS, CONFIGS, etc.) from command line.

\hypertarget{d3_r2q.x-code}{%
\subsection{d3\_r2q.x code}\label{d3_r2q.x-code}}

This code reads the 2nd order FCs and computes a number of different
quantities that only depend on the harmonic 2nd order force constants,
it is currently evolving and should be quite easy to modify and extend
according to your needs. It reads its input from a Fortran namelist
called \&r2qinput which contains the variables listed in the next
section.

After the namelist, the code will look for the keyword QPOINTS and will
start reading the list of q-points. This part is described in detail in
section QPOINTS

\hypertarget{namelist-r2qinput}{%
\subsubsection{Namelist \&r2qinput}\label{namelist-r2qinput}}

\hypertarget{calculation-character-default-freq}{%
\paragraph{\texorpdfstring{calculation (CHARACTER, default:
\enquote{freq})}{calculation (CHARACTER, default: ``freq'')}}\label{calculation-character-default-freq}}

The type of calculation to perform, currently this can be:

\begin{itemize}
\tightlist
\item
  \enquote{freq}: Compute the phonon frequencies, this keyword that you
  specify nq and a QPOINTS section.
\item
  \enquote{jdos}: Compute the joint density of state. You will have to
  specify the integration grid nk and a configuration in the CONFIGS
  section (if you specify more, only the first configuration will be
  used)
\item
  \enquote{rms} : Compute the root mean square displacement of the atoms
  around their equilibrium positions in the harmonic hamiltonian at a
  give temperature. It requires nk and one single configuration in
  CONFIG.
\item
  \enquote{fh}: Compute the phonon free energy or, for T=0, the
  zero-point energy.
\end{itemize}

\hypertarget{prefix-character-default-the-value-of-calculation}{%
\paragraph{prefix (CHARACTER, default: the value of
calculation)}\label{prefix-character-default-the-value-of-calculation}}

The first part of the output file name, the file will be called
\enquote{prefix.out}.

\hypertarget{outdir-character-default-.-i.e.-the-current-directory}{%
\paragraph{\texorpdfstring{outdir (CHARACTER, default: \enquote{./},
I.e. the current
directory)}{outdir (CHARACTER, default: ``./'', I.e. the current directory)}}\label{outdir-character-default-.-i.e.-the-current-directory}}

Location where the output file will be saved.

\hypertarget{file_mat2-character-no-default}{%
\paragraph{file\_mat2 (CHARACTER, no
default)}\label{file_mat2-character-no-default}}

The file of the 2nd order force constants, produced by thermal2 internal
version of q2r.x

\hypertarget{asr2-character-default-no}{%
\paragraph{\texorpdfstring{asr2 (CHARACTER, default:
\enquote{no})}{asr2 (CHARACTER, default: ``no'')}}\label{asr2-character-default-no}}

Method used to apply the acoustic sum rule, can be

\begin{itemize}
\tightlist
\item
  \enquote{no} (do not apply ASR)
\item
  \enquote{simple} (apply the compensation term to the on-site force
  constant)
\end{itemize}

\hypertarget{nq-integer-no-default}{%
\paragraph{nq (INTEGER, no default)}\label{nq-integer-no-default}}

Number of q-points to read (see the QPOINTS section below)

\hypertarget{print_dynmat-logical-default-.false.}{%
\paragraph{print\_dynmat (LOGICAL, default:
.false.)}\label{print_dynmat-logical-default-.false.}}

If set to .true. A file containing the dynamical matrix, in phonon
format, will be saved for each q-point (only works for
calculation=\enquote{freq}). The file name will be
\enquote{r2q\_dyn\_NQ} where NQ is the progressive number of the point.

\hypertarget{sort_freq-character-default-default}{%
\paragraph{\texorpdfstring{sort\_freq (CHARACTER, default:
\enquote{default})}{sort\_freq (CHARACTER, default: ``default'')}}\label{sort_freq-character-default-default}}

When plotting the linewidth and frequencies along a path, there are
several ways to order the frequencies and associated linewidth and
shifted frequencies:

\begin{itemize}
\tightlist
\item
  \enquote{default}: keep the default order of increasing frequencies
\item
  \enquote{overlap}: sort each point in order to maximize the overlap of
  each band polarization with the corresponding band at the previous
  point. This is the best choice for paths, but probably will not work
  for 2D or 3D grid plots
\end{itemize}

\hypertarget{print_velocity-logical-default-.false.}{%
\paragraph{print\_velocity (LOGICAL, default:
.false.)}\label{print_velocity-logical-default-.false.}}

If set to .true. a file containing the phonon group velocities will be
saved (only applies when calculation=\enquote{freq}). The file name will
be prefix\_vel.out, after the path length and the q-point, the
velocities are printed in Cartesian coordinates, Rydberg units
(1.09×10{[}6{]} m·s−1).

\hypertarget{ne-de-e0-sigma_e-integer-real-real-in-cm-1-no-default}{%
\paragraph{ne, de, e0, sigma\_e (INTEGER, REAL, REAL, in cm-1 no
default)}\label{ne-de-e0-sigma_e-integer-real-real-in-cm-1-no-default}}

Used for jdos calculation, see the description in \&lwinput section,
below.

\hypertarget{output-format}{%
\subsubsection{Output format}\label{output-format}}

The r2q.x code will produce an output file for every configuration. The
output files will be named \$prefix.out (where \$prefix is the value of
the input variable prefix). It contains the following columns:

\begin{itemize}
\tightlist
\item
  1 The line number i.e.~point index
\item
  2 The length of the q-point path or, if computing over a grid the
  weight of the q-points
\item
  3→5 The coordinates of the q-point in units of 2π/alat
\item
  6→5+3~nat The phonon frequencies in cm-1.
\end{itemize}

\hypertarget{d3_qha.x-code}{%
\subsection{d3\_qha.x code}\label{d3_qha.x-code}}

This code reads the 2nd order FCs for a series of volumes and computes
the phonon free energy for a given list of temperature, optionally
adding a pV (pressure × volume) hydrostatic term. It then fits the total
free energy with an equation of state to find the equilibrium volume at
each temperature, and find the temperature/volume curve and the
volumetric thermal expansion coefficient.

\hypertarget{namelist-qhainput}{%
\subsubsection{Namelist \&qhainput}\label{namelist-qhainput}}

\hypertarget{calculation-character-default-gibbs}{%
\paragraph{\texorpdfstring{calculation (CHARACTER, default:
\enquote{gibbs})}{calculation (CHARACTER, default: ``gibbs'')}}\label{calculation-character-default-gibbs}}

The type of calculation to perform, at the moment it can only do the
default.

\hypertarget{prefix-character-default-the-value-of-calculation-1}{%
\paragraph{prefix (CHARACTER, default: the value of
calculation)}\label{prefix-character-default-the-value-of-calculation-1}}

The first part of the output file name, the file will be called
\enquote{prefix\ldots.out}. Where the \enquote{\ldots{}} part depends on
the kind of calculation (See the
\protect\hyperlink{output-format-1}{Output format} section)

\hypertarget{outdir-character-default-.-i.e.-the-current-directory-1}{%
\paragraph{\texorpdfstring{outdir (CHARACTER, default: \enquote{./},
i.e.~the current
directory)}{outdir (CHARACTER, default: ``./'', i.e.~the current directory)}}\label{outdir-character-default-.-i.e.-the-current-directory-1}}

Location where the output file will be saved.a

\hypertarget{nt-dt-t0-integer-real-real-in-k}{%
\paragraph{nT, dT, T0 (INTEGER, REAL, REAL, in
K)}\label{nt-dt-t0-integer-real-real-in-k}}

These three variables define the list of temperatures to compute: T0,
T0+dT, \ldots, T0+(nT-1)dT. Use a sufficiently small value for dT in
order to have a reliable thermal expansion coefficient.

\hypertarget{asr2-character-default-no-1}{%
\paragraph{\texorpdfstring{asr2 (CHARACTER, default:
\enquote{no})}{asr2 (CHARACTER, default: ``no'')}}\label{asr2-character-default-no-1}}

Method used to apply the acoustic sum rule, can be \enquote{no} (do not
apply ASR), \enquote{simple} (apply the compensation term to the on-site
force constant). You can use the script tools/apply\_asr.sh in order to
apply more sophisticated sum rules using matdyn.x from qe.

\hypertarget{nk-3x-integer-no-default}{%
\paragraph{nk (3x INTEGER, no default)}\label{nk-3x-integer-no-default}}

The size of the grid used to integrate the phonon-phonon interaction
processes.

\hypertarget{grid_type-character-default-simple}{%
\paragraph{\texorpdfstring{grid\_type (CHARACTER, default:
\enquote{simple})}{grid\_type (CHARACTER, default: ``simple'')}}\label{grid_type-character-default-simple}}

Set this to \enquote{simple} to use a regular unshifted grid in
reciprocal space. See the
\protect\hyperlink{grid_type-character-default-simple-1}{description of
d3\_lw.x input} for more details about this option.

\hypertarget{press_kbar-real-kbar-0}{%
\paragraph{press\_kbar (REAL, kbar, 0)}\label{press_kbar-real-kbar-0}}

\hypertarget{press_gpa-real-gpa-0}{%
\paragraph{press\_Gpa (REAL, Gpa, 0)}\label{press_gpa-real-gpa-0}}

Optionally add an hydrostatic pressure, which will contribute a term pV
to the total energy. The sign convention for pressure is that higher
positive pressure means pushing stronger on the sample (i.e.~you may
want to use a positive pressure value 99\% of the times).

\hypertarget{n_volumes-integer-no-default}{%
\paragraph{n\_volumes (INTEGER, no
default)}\label{n_volumes-integer-no-default}}

The number of volumes that have been computed ab-initio. The code expect
to find a list of n\_volumes force constant files and total electronic
energies after the namelist.

\hypertarget{eos-character-default-murn}{%
\subsubsection{\texorpdfstring{eos (CHARACTER, default:
\enquote{murn})}{eos (CHARACTER, default: ``murn'')}}\label{eos-character-default-murn}}

Kind of equation of state to use:

\begin{enumerate}
\def\labelenumi{\arabic{enumi}.}
\tightlist
\item
  murn → Murnaghan
\item
  birch1 → Birch 1st order
\item
  birch3 → Birch 3rd order
\item
  keane → Keane
\end{enumerate}

Note that in principle the EOS are empiriclly suitable for V-P
(volume-pressure) curves, not for V-T. In practice they work remarkably
well in all the cases we have tested.

\hypertarget{list-of-volumes}{%
\subsubsection{List of volumes}\label{list-of-volumes}}

After the namelist, the code reads n\_volumes lines, each line contains
the name of a force-constant file (produced by d3\_q2r.x) and the total
electronic energy corresponding to it. For examples:

\begin{verbatim}
&qhainput
 ...
 n_volumes = 3
/
mat2R_1   -85.72256763
mat2R_2   -85.72385873
mat2R_3   -85.73049838
\end{verbatim}

It may be necessary to enclose the name of the file in quotes
\enquote{\ldots{}} if it contains any special character, such as
\enquote{/}. The energy are in Ry and are just the \enquote{total
energy} printed by pw.x at total convergence.

\hypertarget{output-format-1}{%
\subsubsection{Output format}\label{output-format-1}}

The code will create, in outdir, a file for each temperature with the
equation of state for that temperature, and a final file with the
theoretical volume/temperature curve, For each temperature, a file
called \(prefix_T\)temperature.dat, in the file header you will find
these informations that are obtained by fitting the total Gibbs free
energy with an equation of state:

\begin{itemize}
\tightlist
\item
  v0: the equilibrium volume at this temperature
\item
  g0: the theoretical minimu of the free energy at v0
\item
  k0, dk0, d2k0: bulk modulus, its first and second derivatives
\end{itemize}

Afterwards, the file will contain these columns: 1. Cell volume in bohr
2. total Gibbs temperature dependent energy (g) 3. total electronic
energy (from input) 4. zero-point vibrational energy 5. phonon free
energy 6. hydrostatic contribution (pV) 7. total energy w.r.t the
minimum of the equation of state fitted (g-g0) 8. total Gibbs free
energy from the fit

The final file, called \$prefix.dat, contains the following columns: 1.
Temperature\\
2. Volume (bohr3)\\
3. Volumetric thermal expansion\\
4. Theoretical total free energy, at the equilibrium volume for this
temperature (g(T,v0)) 5. Bulk modulus, its first and second derivatives

\hypertarget{d3_lw.x-code}{%
\subsection{d3\_lw.x code}\label{d3_lw.x-code}}

This code can compute the intrinsic phonon-phonon interaction and the
interaction of phonons with isotopic disorder and border scattering. It
reads its input variables from the \&lwinput namelist and from the
QPOINTS, CONFIGS and ISOTOPES lists, which are described in detail in
their corresponding sections.

\hypertarget{namelist-lwinput}{%
\subsubsection{Namelist \&lwinput}\label{namelist-lwinput}}

\hypertarget{calculation-character-default-lw}{%
\paragraph{\texorpdfstring{calculation (CHARACTER, default:
\enquote{lw})}{calculation (CHARACTER, default: ``lw'')}}\label{calculation-character-default-lw}}

The type of calculation to perform, it can take several different
values:

\begin{itemize}
\tightlist
\item
  \enquote{lw imag}: compute the imaginary part of the self-energy,
  i.e.~the phonon linewidth. In this mode the code will enforce
  conservation of energy with a Diract delta function approximated with
  a Gaussian function of width delta, read from the CONFIGS section.
\item
  \enquote{lw full}: compute the entire self energy, the real part is
  the lineshift and the imaginary part is the linewidth. In this case
  the value of delta from the CONFIGS will be used as a regularization
  for the self-energy.
\item
  \enquote{spf full}: compute the spectra function, also known as σ(ω),
  for a list of energies, provided by the input variables ne, de, e0 and
  sigma\_e (see below) and for all q-points.
\item
  \enquote{spf imag}: as \enquote{spf full}, but only the imaginary part
  of the self-energy will be used, i.e.~the spectra function will be
  centered around the non-shifted phonon energy. This is often in better
  agreement with experiments than \enquote{spf full}, unless you also
  include somehow the 4-phonons self-energy contribution, because the
  real part of the 3-phonon term and the 4-phonon terms tend to cancel
  each other out.
\item
  \enquote{spf simple}: simulate the spectral function as a
  superposition of Lorentzian functions centered around the phonon
  frequencies and appropriate width, see also ne, de and n0. This is
  equivalent to \enquote{spf full} when the anharmonicity is very weak.
\item
  \enquote{final}: decompose the contribution to the linewidth to a
  specific energy and q-point (specified with the e\_initial and
  q\_initial keywords) over the energy of the final states in the
  scattering process or over the final q, or both (see q\_summed and
  q\_resolved). In the first case you must specify the range of final
  energies to consider with the ne, de and e0 keywords; in the second
  case the possible final q-points will be read from the QPOINTS
  section.
\end{itemize}

\hypertarget{prefix-character-default-the-value-of-calculation-2}{%
\paragraph{prefix (CHARACTER, default: the value of
calculation)}\label{prefix-character-default-the-value-of-calculation-2}}

The first part of the output file name, the file will be called
\enquote{prefix\ldots.out}. Where the \enquote{\ldots{}} part depends on
the kind of calculation (See the
\protect\hyperlink{output-format-2}{Output format section})

\hypertarget{outdir-character-default-.-i.e.-the-current-directory-2}{%
\paragraph{\texorpdfstring{outdir (CHARACTER, default: \enquote{./},
i.e.~the current
directory)}{outdir (CHARACTER, default: ``./'', i.e.~the current directory)}}\label{outdir-character-default-.-i.e.-the-current-directory-2}}

Location where the output file will be saved.

\hypertarget{file_mat2-character-no-default-1}{%
\paragraph{file\_mat2 (CHARACTER, no
default)}\label{file_mat2-character-no-default-1}}

The file of the 2nd order force constants, produced by thermal2 internal
version of q2r.x

\hypertarget{file_mat3-character-no-default}{%
\paragraph{file\_mat3 (CHARACTER, no
default)}\label{file_mat3-character-no-default}}

The file of the 3rd order force constants, produced by qq2rr.x or asr3.x
or sparse.x

\hypertarget{asr2-character-default-no-2}{%
\paragraph{\texorpdfstring{asr2 (CHARACTER, default:
\enquote{no})}{asr2 (CHARACTER, default: ``no'')}}\label{asr2-character-default-no-2}}

Method used to apply the acoustic sum rule, can be \enquote{no} (do not
apply ASR), \enquote{simple} (apply the compensation term to the on-site
force constant).

\hypertarget{nq-integer-no-default-1}{%
\paragraph{nq (INTEGER, no default)}\label{nq-integer-no-default-1}}

Number of q-points to read (see below)

\hypertarget{nconf-integer-no-default}{%
\paragraph{nconf (INTEGER, no default)}\label{nconf-integer-no-default}}

Number of configurations to read in the CONFIGS section, see the
detailed description below.

\hypertarget{nk-3x-integer-no-default-1}{%
\paragraph{nk (3x INTEGER, no
default)}\label{nk-3x-integer-no-default-1}}

The size of the grid used to integrate the phonon-phonon interaction
processes.

\hypertarget{grid_type-character-default-simple-1}{%
\paragraph{\texorpdfstring{grid\_type (CHARACTER, default:
\enquote{simple})}{grid\_type (CHARACTER, default: ``simple'')}}\label{grid_type-character-default-simple-1}}

Set this to \enquote{simple} to use a regular unshifted grid in
reciprocal space.

Use \enquote{random} to use a grid shifted by a random vector; the
random shift is not applied to the directions where there is only one
k-point, i.e.~if you have a N×M×1 grid, there will be no shift along z.
Using a random shifted grid, can easily reduce the number of points
required for convergence by half in each direction, for a speed-up of 8
for linewidth calculations and 64 times for tk calculations; however, it
will break symmetry, and give (slightly) different thermal conductivity
for different directions, even in highly symmetric crystals. See also
\protect\hyperlink{xk0-3x-real-default-000}{xk0} to manually apply a
grid shift. Shifted grids are currently disabled for CGP calculations,
the reason is that we did not manage to impose the detailed balance
condition in this case, causing convergency issues and runaway
minimization. If you are at the limit of computational power, it is
possible to disable this limitation in the code, but a careful
examination of the minimization procedure must be done.

Set to \enquote{bz} to use a grid centered in the Brillouin zone. This
option will duplicate the points that are on the boundary of the BZ and
assign them an appropriate integration weight. Using a BZ grid should
eliminate a possible source of unwanted symmetry breaking, although it
adds some complexity, uses more points, and its effectiveness is not
evident in practice. It can be useful for doing plots, but apart from
this case, we recommend using \enquote{simple} instead.

\hypertarget{xk0-3x-real-default-000}{%
\paragraph{xk0 (3x REAL, default:
0,0,0)}\label{xk0-3x-real-default-000}}

A shift to apply to the grid of k-points, in units of half the lattice
spacing; i.e.~set it to (1,1,1) to have the standard Monkhorst-Pack
shifted grid, any fractional values is allowed.

\hypertarget{ne-de-e0-sigma_e-integer-real-real-in-cm-1-no-default-1}{%
\paragraph{ne, de, e0, sigma\_e (INTEGER, REAL, REAL, in cm-1 no
default)}\label{ne-de-e0-sigma_e-integer-real-real-in-cm-1-no-default-1}}

When doing a Spectral function (d3\_lw.x) or Final state (d3\_lw.x ) or
Joint-DOS (d3\_r2q.x) calculation you have to define an energy axis with
these variables. The axis will include ne equally spaced points starting
from e0 and up to e0+(ne-1)de. The sum over the q-points will be
convoluted with a gaussian of width sigma\_e (default: 5~de) to obtain a
smooth curve.

\hypertarget{nu_initial-1..3nat-no-default}{%
\paragraph{nu\_initial (1..3×nat, no
default)}\label{nu_initial-1..3nat-no-default}}

When doing a Final state decomposition calculation, this is the band of
the initial state considered in the scattering process. If not
specified, the final state will be computed for all initial bands at the
energ specified by e\_initial, but as a consequence it will not be
decomposed over the final bands.

\hypertarget{e_initial-in-cm-1-no-default}{%
\paragraph{e\_initial (in cm-1, no
default)}\label{e_initial-in-cm-1-no-default}}

When doing a Final state decomposition calculation, this is the energy
of the initial state considered in the scattering process. If not
specified, the energy of the phonon corresponding to nu\_initial will be
used.

\hypertarget{q_initial-3x-real-in-units-of-2ux3c0alat}{%
\paragraph{q\_initial (3x REAL, in units of
2π/alat)}\label{q_initial-3x-real-in-units-of-2ux3c0alat}}

As e\_initial, specifies the initial q-point.

\hypertarget{q_summed-logical-default-false}{%
\paragraph{q\_summed (LOGICAL, default:
false)}\label{q_summed-logical-default-false}}

When doing a final state calculation, set to true to project the
infinitesimal contribution to the linewidth over the given final
q-points, the energy dependence is integrated out. The analysis is
performed over the q-points specified in the QPOINTS section, you can
use either a path or a grid, depending on the kind of plot you want,
using a xsf or bxsf kind of grid can be used to do a 3D plot with
\href{http://www.xcrysden.org}{XCrysDen}. The full grid, specified by
nk, will still be used to compute the linewidth, but te contribution
will be projected to each q-point in the list with a Gaussian smearing
given by sigmaq.

Note that even high-symmetry points can, and often do, decay toward
lower symmetry points; a high-symmetry path can easily miss the most
important final states. We recommend using a relatively coarse grid
first, you will be able to spot the most favored points by sorting the
output file. I.e. this command: sort -gk 5 final\_T300\_s1.out will
output as the final lines the most important decay processes. See also
sigmaq and q\_resolved and the output format. See also the
\protect\hyperlink{compute-the-spectral-function-along-a-path}{examples}.

\hypertarget{q_resolved-logical-default-false}{%
\paragraph{q\_resolved (LOGICAL, default:
false)}\label{q_resolved-logical-default-false}}

As q\_summed, but the energy dependence is not integrated out. The
output file will contain an analysis of the decay process as a function
of the final state q-point and energy. Note that this file can be huge,
and it is almost impossible to plot for an entire grid. Along a line,
you can produce a color plot using this gnuplot command (assuming that
the file freq.out contains the frequencies):

\begin{verbatim}
set palette defined (0  "white", 1  "red") 
plot "fs_qresolved_T300_s1.out" u 1:2:6 w image, \
     for [i=6:11] 'freq.out' u 2:i w l lt -1 not
\end{verbatim}

See also sigmaq and q\_summed and the output format. Note that the
\enquote{image} plot mode suppose that the x-axis spacing is constant,
if it is not, you will have to do a 3D \enquote{splot} with
\enquote{view map} in order to obtain a good plot.

\hypertarget{sigmaq-real-default-0.1-in-units-of-2ux3c0alat}{%
\paragraph{sigmaq (REAL, default: 0.1 in units of
2π/alat)}\label{sigmaq-real-default-0.1-in-units-of-2ux3c0alat}}

Used in conjunction with q\_summed or q\_resolved to obtain a nice
smooth plot over the q-points by convoluting it with a Gaussian function
of width sigmaq. Set it to something of the order of the spacing between
the q-points.

\hypertarget{exp_t_factor-logical-default-false}{%
\paragraph{exp\_t\_factor (LOGICAL, default:
false)}\label{exp_t_factor-logical-default-false}}

UNTESTED/EXPERIMENTAL! When doing a spectra function calculation, add an
elastic peak of equation (1 + f\_bose(e,T)) / e) which should emulate
the elastic peak of neutron spectroscopy.

\hypertarget{sort_freq-character-default-default-1}{%
\paragraph{\texorpdfstring{sort\_freq (CHARACTER, default:
\enquote{default})}{sort\_freq (CHARACTER, default: ``default'')}}\label{sort_freq-character-default-default-1}}

When plotting the linewidth and frequencies along a path, there are
several ways to order the frequencies and associated linewidth and
shifted frequencies:

\begin{itemize}
\tightlist
\item
  \enquote{default}: keep the default order of increasing frequencies
\item
  \enquote{overlap}: sort each point in order to maximize the overlap of
  each band polarization with the corresponding band at the previous
  point. This is the best choice for paths, but probably will not work
  for 2D or 3D grid plots
\item
  \enquote{shifted}: sort in order of shifted frequencies,
  i.e.~frequency+lineshift. It is only meaningful when doing a
  \enquote{lw full} calculation, it can help to have good quality plots
  when \enquote{overlap} fails.
\end{itemize}

\hypertarget{isotopic_disorder-logical-default-false}{%
\paragraph{isotopic\_disorder (LOGICAL, default:
false)}\label{isotopic_disorder-logical-default-false}}

Set this to true to include scattering from isotopic disorder. You will
also have to specify the isotopic composition of every element in the
system in the ISOTOPES section. NOTE: isotopes are only used for
linewidth calculation, they are no used for spectral functions and final
state decomposition.

\hypertarget{casimir_scattering-logical-default-false}{%
\paragraph{casimir\_scattering (LOGICAL, default:
false)}\label{casimir_scattering-logical-default-false}}

Set this variable to true to include scattering with boundary, treated
with the Casimir formula. See the following variables for detail on how
to specify the boundary structure. Casimi scattering is only applied to
linewidth calculation, it has no effect on spectra function and final
state calculations.

\hypertarget{sample_dir-3x-real-normalized-default-no-direction}{%
\paragraph{sample\_dir (3x REAL, normalized, default: no
direction)}\label{sample_dir-3x-real-normalized-default-no-direction}}

The direction in which Casimir scattering is prevalent. Set it to zero
to consider omnidirectional scattering, i.e.~scattering from grains, or
to a specific direction to consider a wire. More complex geometries are
not implemented yet, please let us know if you are interested.

\hypertarget{sample_length_au-real-bohr-no-default}{%
\paragraph{sample\_length\_au (REAL, bohr, no
default)}\label{sample_length_au-real-bohr-no-default}}

\hypertarget{sample_length_mu-real-micrometers-no-default}{%
\paragraph{sample\_length\_mu (REAL, micrometers, no
default)}\label{sample_length_mu-real-micrometers-no-default}}

\hypertarget{sample_length_mm-real-millimeters-no-default}{%
\paragraph{sample\_length\_mm (REAL, millimeters, no
default)}\label{sample_length_mm-real-millimeters-no-default}}

The average scattering length in the Casimir model; you can specify only
one of the three variables, according to the unit of measure you prefer.
Note that the Caimir model also include a structure factor, usually set
to 2, which is not included in our implementation; you will have to
include it directly in the gain size in input.

\hypertarget{max_seconds-no-default-in-seconds}{%
\paragraph{max\_seconds (no default, in
seconds)}\label{max_seconds-no-default-in-seconds}}

\hypertarget{max_time-format-hh.mmss}{%
\paragraph{max\_time (format hh.mmss)}\label{max_time-format-hh.mmss}}

The maximum running time after which the code will stop, you can only
set one of the two. Notice that only tk.x includes a restart mechanism.

\hypertarget{output-format-2}{%
\subsubsection{Output format}\label{output-format-2}}

The lw.x code will produce an output file for every configuration. The
output files will be named \$prefix\_T\$XX\_s\$YY.out Where \$XX is the
temperature in Kelvin and \$YY the delta in cm-1. Each output file will
contain a relatively large number of columns, depending on the
calculation.

NOTE: In all of the following cases case when specifying QPOINTS as
\enquote{grid} or \enquote{bz}, the length of the path will actually be
replaced by the weight of the point used to do an integral in reciprocal
space. Calculation \enquote{lw imag} Number of the column, or columns
and its and content:

\begin{itemize}
\tightlist
\item
  1 The line number i.e.~point index
\item
  2 The length of the q-point path (when doing a path) or the weight of
  the q-point (when doing a grid calculation)
\item
  3→5 The coordinates of the q-point in 2π/alat.
\item
  6→5+3~nat The unperturbed phonon frequencies in cm-1
\item
  6+3~nat→5+6~nat The linewidth (HWHM) in cm-1
\end{itemize}

Calculation \enquote{lw full}:

\begin{itemize}
\tightlist
\item
  1→5+6~nat. Same content as \enquote{lw imag}
\item
  6+6~nat→5+8~nat.
\end{itemize}

The shifted phonon frequencies (i.e.~frequency+shift) in cm-1. If the
sort\_shifted\_q keyword was set to true, the shifted frequencies are
sorted in increasing order and the corresponding linewidth are sorted
accordingly. The unperturbed frequencies are left unchanged.

Calculation \enquote{spf}:

\begin{itemize}
\tightlist
\item
  1 The energy axis in cm-1
\item
  2 The length of the q-point path (when doing a path) or the weight of
  the q-point (when doing a grid calculation)
\item
  3 The total spectral function in 1/cm-1
\item
  4→3+3~nat Contribution to the spectral function from each band in
  1/cm-1
\end{itemize}

Note that the energy axis cycles faster than the path length and there
is a whitespace after each q-point which makes plotting with gnuplot
\enquote{pm3d} style easy with this syntax (column 3 is used twice, both
for heigth and colour): sp \enquote{file} u 1:2:3:3 w pm3d

Calculation \enquote{final}, if nu\_initial \textbf{is not} specified:

\begin{itemize}
\tightlist
\item
  1 The energy axis in cm-1
\item
  2 The total final state weight in 1/cm-1
\item
  3 Cohalescence part of the total final state weight in 1/cm-1
\item
  4 Scattering part of the total final state weight in 1/cm-1
\item
  5→4+3 nat Contribution to the final state \textbf{from} each band in
  1/cm-1
\item
  5+3nat→4+6 nat Cohalescence to the final state \textbf{from} each band
  in 1/cm-1
\item
  5+6nat→4+9 nat Scattering to the final state \textbf{from} each band
  in 1/cm-1
\end{itemize}

Calculation \enquote{final}, if nu\_initial \textbf{is} specified:

\begin{itemize}
\tightlist
\item
  1 The energy axis in cm-1
\item
  2 The total final state weight in 1/cm-1
\item
  3 Cohalescence part of the total final state weight in 1/cm-1
\item
  4 Scattering part of the total final state weight in 1/cm-1
\item
  5→4+3 nat Contribution to the final state \textbf{to} each band in
  1/cm-1
\item
  5+3nat→4+6 nat Cohalescence to the final state \textbf{to} each band
  in 1/cm-1
\item
  5+6nat→4+9 nat Scattering to the final state \textbf{to} each band in
  1/cm-1
\end{itemize}

Calculation \enquote{final}, q\_resolved TRUE:

\begin{itemize}
\tightlist
\item
  1 The energy axis in cm-1
\item
  2 The length of the q-point path (when doing a path, units 2π/alat) or
  the weight of the q-point (when doing a grid calculation)
\item
  3→5 The coordinates of the q-point in (2π/alat).
\item
  6 The infinitesimal contribution to the linewidth, as a function of
  energy and q (1/cm-1)
\item
  7→6+3~nat The infinitesimal contribution to the linewidth, decomposed
  by band (1/cm-1)
\end{itemize}

Calculation \enquote{final}, q\_summed TRUE:

\begin{itemize}
\tightlist
\item
  1 The length of the q-point path (when doing a path, units 2π/alat) or
  the weight of the q-point (when doing a grid calculation)
\item
  2→4 The coordinates of the q-point in (2π/alat).
\item
  5 The infinitesimal contribution to the linewidth from this q-point
  (1/cm-1)
\item
  6→5+3~nat The infinitesimal contribution to the linewidth, decomposed
  by band (1/cm-1)
\end{itemize}

\hypertarget{d3_tk.x-code}{%
\subsection{d3\_tk.x code}\label{d3_tk.x-code}}

The d3\_tk.x code can compute the thermal conductivity coefficient in
the SMA or by exact diagonalization of the BTE. Most of its input
variables are the same as d3\_lw.x, we refer to the previous section for
their description

\hypertarget{namelist-tkinput}{%
\subsubsection{Namelist \&tkinput}\label{namelist-tkinput}}

Most of the variable used by d3\_tk.x are also used for d3\_lw.x. The
following variables have the same meaning as in d3\_lw.x: outdir,
prefix, file\_mat2, file\_mat3, asr2, nconf, nk, grid\_type, xk0,
isotopic\_disorder, casimir\_scattering, sample\_direction,
sample\_length\_* The following variables are specific of d3\_tk.x.

\hypertarget{calculation-character-default-sma}{%
\paragraph{\texorpdfstring{calculation (CHARACTER, default:
\enquote{sma})}{calculation (CHARACTER, default: ``sma'')}}\label{calculation-character-default-sma}}

Set this variable to \enquote{sma} to compute use the single mode
approximation or to \enquote{cgp} to do an iterative diagonalization of
the BTE with the Conjugate Gradient algorithm with preconditioning.
Please note that \enquote{cgp} algorithm will also output the SMA
thermal conductivity at its first iteration, but computed in a slightly
different way which is slower but guarantees the phonon-phonon
scattering matrix to be well defined and the functional to be minimized
to be positive definite.

\hypertarget{nk_in-3x-integer-default-same-value-as-nk}{%
\paragraph{nk\_in (3x INTEGER, default: same value as
nk)}\label{nk_in-3x-integer-default-same-value-as-nk}}

When doing a \enquote{sma} calculation, the linewidth will be computed,
for each point in the nk grid, integrating over the nk\_in grid. Thermal
conductivity will then be integrated over the nk grid. In principle,
there is no reason for the two grids to be identical, as one quantity
could be harder to converge than the other. When doing a CGP
calculation, the inner and outer grids must be identical, nk\_in will be
ignored and nk used instead.

\hypertarget{grid_type_in-character-default-same-value-as-grid_type}{%
\paragraph{grid\_type\_in (CHARACTER, default: same value as
grid\_type)}\label{grid_type_in-character-default-same-value-as-grid_type}}

Same as grid\_type, but applied to the inner grid. Ignored for
\enquote{cgp} calculations.

\hypertarget{xk0_in-3x-real-default-same-value-as-xk0}{%
\paragraph{xk0\_in (3x REAL, default: same value as
xk0)}\label{xk0_in-3x-real-default-same-value-as-xk0}}

Same as xk0, but applied to the inner grid. Ignored for \enquote{cgp}
calculations.

\hypertarget{thr_tk-real-default-1.d-4-wmk}{%
\paragraph{thr\_tk (REAL, default: 1.d-4
W/mK)}\label{thr_tk-real-default-1.d-4-wmk}}

Threshold on the convergence of thermal conductivity the of the CGP
minimization when solving the BTE. The default value should be enough to
get the thermal conductivity with 4 or 5 significant digits.

\hypertarget{niter_max-integer-default-1000}{%
\paragraph{niter\_max (INTEGER, default:
1000)}\label{niter_max-integer-default-1000}}

Maximum number of iteration of the CGP algorithm. Normally less than 100
iterations are needed.

\hypertarget{store_lw-logical-default-.false.}{%
\paragraph{store\_lw (LOGICAL, default:
.false.)}\label{store_lw-logical-default-.false.}}

Set this to true, when doing a SMA calculation, and the code will write
to disk the values of the intrinsic, isotopi and casimir linewidths at
the end of the calculation. In the tool/ you can find an octave/matlab
script recompute\_sma.m to recompute the SMA tk from these files. This
can be used to inexpensively test different isotopic and Casimir
scattering parameters without repeating the expensive intrinsic
linewidth calculation. See also intrinsic\_scattering

\hypertarget{intrinsic_scattering-logical-default-.true.}{%
\paragraph{intrinsic\_scattering (LOGICAL, default:
.true.)}\label{intrinsic_scattering-logical-default-.true.}}

If this parameter is set to false, during a SMA calculation, the code
will skip the calculation of the intrinsic phonon-phonon scattering. At
the end, thermal conductivity cannot not be computed, hence be sure to
set store\_lw=.true., or the calculation will be wasted, in order to
store isotopic and Casimir linewidths.

\hypertarget{restart-logical-default-.false.}{%
\paragraph{restart (LOGICAL, default:
.false.)}\label{restart-logical-default-.false.}}

Set this to true, when doing a CGP calculation, and the code will write
to disk the state of the minimization at each iteration, it will then be
able to restart from the last step. Set this to .false. will have two
effects: 1) the code will ignore any restart information already present
on disk and 2) the code will not write restart information on disk.
Restart information can take quite a bit of space for dense grids.

\hypertarget{mfp_cutoff-logical-default-false}{%
\paragraph{mfp\_cutoff (LOGICAL, default:
false)}\label{mfp_cutoff-logical-default-false}}

Set this variable to true to include scattering with boundary, by
cutting off the contribution of all phonon modes with a mean free path
(MFP, inverse full linewidth times group velocity) larger than the
sample size. This approach is simpler than using casimir\_scattering,
but it assumes that all phonons with a MFP longer than the sample will
scatter, and all those shorter will not; i.e.~that all phonons with a
MFP longer than the sample hav an infinite linewidth; this is only true
in the ballistic regime, which is not the aim of this software anyway.
We recommend using casimir\_scattering instead. This option will use the
input variable sample\_dir and sample\_length\_\# to determine the
cutoff length. This option only works for SMA thermal conductivity.

\hypertarget{volume_factor-real-default-1}{%
\paragraph{volume\_factor (REAL, default:
1)}\label{volume_factor-real-default-1}}

A dimensionless parameter to rescale the volume of the crystal unit
cell. When studying a 2D material, it is useful to normalize the thermal
conductivity with the volume of bulk, excluding the vacuum space left
between periodic copies of the 2D slab. I.e. if the bulk material has an
inter-layer spacing of H and you have built your 2D slab geometry with a
vacuum distance V, you have to set volume\_factor=H/V.

\hypertarget{output-format-3}{%
\subsubsection{Output format}\label{output-format-3}}

\hypertarget{sma-calculation}{%
\paragraph{SMA calculation}\label{sma-calculation}}

When doing a SMA calculation the tk.x code will produce four output
files:

\begin{enumerate}
\def\labelenumi{\arabic{enumi}.}
\tightlist
\item
  A file named \(prefix.\)grid\_size\_TOT.out, where \$prefix is the
  input value of prefix and \$grid\_size is the size of the integration
  grid (e.g.~\enquote{20x20x20}). This file will contain one line per
  configuration, in each line you will:
\end{enumerate}

\begin{itemize}
\tightlist
\item
  1 the configuration number,
\item
  2 the value of sigma
\item
  3 the temperature
\item
  4→6 the diagonal elements of the thermal conductivity Kxx, Kyy and Kzz
\item
  7→12 The off-diagonal elements of K, in this order: Kxy, Kxz, Kyz ,
  Kyx, Kzx, Kzy.
\end{itemize}

Two more files with \_P\_sma and \_C instead of \_TOT will contain the
Phonon contribution (dominant in crystals) and Coherent contribution
(dominant in glasses) respectively. See Reference
\protect\hyperlink{ref18}{18} for details.

If the option store\_lw is used, several more, potentially very large,
files will be created. They contain all the quantities required to
recompute the SMA thermal conductivity:

\begin{itemize}
\tightlist
\item
  1 q.\(prefix.\)grid\_size.out: the list of q-vectors (3 columns, in
  cartesian coordinates of 2π/alat) and their respective weight (1
  columns)
\item
  2 freq.\(prefix.\)grid\_size.out: the phonon frequencies (in cm-1, 3x
  number of atoms columns)
\item
  3 lw.\(prefix.\)grid\_size.out: the phonon FWHM (in cm-1, 3x number of
  atoms columns)
\item
  4 vel.\(prefix.\)grid\_size.out: the phonon groups velocity, x,y and z
  for each band (in Rydberg units, 9x number of atoms columns)
\end{itemize}

In the \enquote{tools} directory you can find a mathlab/octave script
\protect\hyperlink{recompute-smam}{recompute\_sma.m} to inexpensively
recompute the thermal conductivity starting from these files.

\hypertarget{cgp-calculation}{%
\paragraph{CGP calculation}\label{cgp-calculation}}

When a CGP calculation several files are created: one with the results
at the last iterations for all the configurations and one file for each
configurations with the results at each iteration. The thermal
conductivity K is always in W/(m·K).

\begin{enumerate}
\def\labelenumi{\arabic{enumi}.}
\tightlist
\item
  A file named \(prefix.\)grid\_size\_TOT.out, where \$prefix is the
  input value of prefix and \$grid\_size is the size of the integration
  grid (e.g.~\enquote{20x20x20}). This file will contain the results
  from the last completed iteration of the code, one line per
  configuration, with these columns:
\end{enumerate}

\begin{itemize}
\tightlist
\item
  1 the configuration number,
\item
  2 the value of sigma
\item
  3 the temperature
\item
  4-6 the diagonal elements of the thermal conductivity Kxx, Kyy and Kzz
\item
  7-12 the off-diagonal elements of K, in this order: Kxy, Kyz, Kyz ,
  Kyx, Kzx, Kzy.
\end{itemize}

\begin{enumerate}
\def\labelenumi{\arabic{enumi}.}
\setcounter{enumi}{1}
\tightlist
\item
  Two more files with \_P\_sma and \_C instead of \_TOT contain the
  Phonon at the SMA level (i.e.~at the first iteration, dominant in
  crystals) and Coherent (dominant in harmonic glasses) contribution to
  thermal conductivity.
\item
  A file for every input configuration, named
  \(prefix.\)grid\_size\_s\(XX_T\)YY.out, where \$XX is the smearing in
  cm-1 and \$YY is the temperature in Kelvin. A line is appended to each
  file at each iteration. The columns are the same as the previous file,
  except that column 1 contains the iteration number.
\end{enumerate}

\hypertarget{d3_tdph.x-code}{%
\subsection{d3\_tdph.x code}\label{d3_tdph.x-code}}

This code reads a set of initial dynamical matrices for a given system
and optimizes the harmonic force constants over a series of images that
can be the output of a molecular dynamics calculation performed with
Quantum ESPRESSO, or of a Langevin Dynamics calculation from the PIOUD
code. The code will expect that the simulation supercell is the same for
the force constants and the dynamics simulations.

\hypertarget{namelist-tdphinput}{%
\subsubsection{Namelist \&tdphinput}\label{namelist-tdphinput}}

\hypertarget{ai-character-default-md}{%
\paragraph{\texorpdfstring{ai (CHARACTER, default
\enquote{md})}{ai (CHARACTER, default ``md'')}}\label{ai-character-default-md}}

Select if the samplig comes from a QE molecular dynamics run (in this
case, read the output of QR from fmd) or from a PIOUD calculation (reads
from files fforce, ftau and ftoten)

\hypertarget{fforce-ftau-ftoten-character-default-forces.dat-positions.dat-pioud.dat}{%
\paragraph{\texorpdfstring{fforce, ftau, ftoten (CHARACTER, default
\enquote{forces.dat}, \enquote{positions.dat},
\enquote{pioud.dat})}{fforce, ftau, ftoten (CHARACTER, default ``forces.dat'', ``positions.dat'', ``pioud.dat'')}}\label{fforce-ftau-ftoten-character-default-forces.dat-positions.dat-pioud.dat}}

Output files from PIOUD containing the ab-initio forces, ions
coordinates and total energy.

\hypertarget{fmd-character-default-md.out}{%
\paragraph{\texorpdfstring{fmd (CHARACTER, default
\enquote{md.out})}{fmd (CHARACTER, default ``md.out'')}}\label{fmd-character-default-md.out}}

Output file from QE

\hypertarget{file_mat2-character-default-mat2r}{%
\paragraph{\texorpdfstring{file\_mat2 (CHARACTER, default
\enquote{mat2R})}{file\_mat2 (CHARACTER, default ``mat2R'')}}\label{file_mat2-character-default-mat2r}}

File of the force constants. Must be in \enquote{periodic} form
(i.e.~generated using -f 0 with d3\_q2r.x)

\hypertarget{nfirst-nskip-nread-integer-default-1-100-50000}{%
\paragraph{nfirst, nskip, nread (INTEGER, default 1, 100,
50000)}\label{nfirst-nskip-nread-integer-default-1-100-50000}}

When reading the MD or PIOUD file, read one every nskip steps starting
from nfirst until nread steps are read in total.

\hypertarget{nmax-integer-default-50000}{%
\paragraph{nmax (INTEGER, default:
50000)}\label{nmax-integer-default-50000}}

Instead of specifying nread, one can use nmax to stop reading after
reachin step nmax of the MD simulation.

\hypertarget{e0-real-default-0}{%
\subsubsection{e0 (REAL, default: 0)}\label{e0-real-default-0}}

A constant to remove from the total energy, i.e.~the ground state
energy. Useful for plotting but has no direct effect on the results.

\hypertarget{basis-character-default-mu}{%
\subsubsection{\texorpdfstring{basis (CHARACTER, default
\enquote{mu})}{basis (CHARACTER, default ``mu'')}}\label{basis-character-default-mu}}

When building the basis of symmetric dynamical matrices, ouse one of
these initial guesses:

\begin{itemize}
\tightlist
\item
  \enquote{simple} orthogonal matrices hermitean matrices with a single
  diagonal or two non-diagonal non-zero elements
\item
  \enquote{mu} from the eignevectors of the original dynamical matrices,
  from file\_mat2, discard the acoustic modes at q=0.
\item
  \enquote{random} start from random matricex (rarely works because of
  accidental degeneracies)
\end{itemize}

\hypertarget{minimization-character-default-ph}{%
\subsubsection{\texorpdfstring{minimization (CHARACTER, default
\enquote{ph})}{minimization (CHARACTER, default ``ph'')}}\label{minimization-character-default-ph}}

When minimizing the phonon parameters, use one of these methods:

\begin{itemize}
\tightlist
\item
  \enquote{ph} only minimize the phonon degrees of freedom
\item
  \enquote{ph+zstar} first minimize the phonons, then the effective
  charges degrees of freedom. It gives rarely any improvement over
  \enquote{ph} because the efefctive charges do not depend strongly on
  temperature.
\item
  \enquote{global} minimize all degrees of freedom simultaneously, this
  method can (and often will) produce unphysical solution if the
  supercell is not huge.
\end{itemize}

\hypertarget{randomization-real-default-0.}{%
\subsubsection{randomization (REAL, default
0.)}\label{randomization-real-default-0.}}

Add some random noise to the initial phonon parameters. If it is a
positive number, all parameters will be randomized with the same
amplitude (proportional to its value times the largest parameter). If it
is a negative value, each parameter will be multiplied by a random
factor between -randomizaiton and +randomization.

\hypertarget{output-format-4}{%
\subsection{Output format}\label{output-format-4}}

The code will produce two force constant files:

\begin{itemize}
\tightlist
\item
  \enquote{matOUT.center} in centered format, suitable for Fourier
  interpolation, band dispersion, etc.
\item
  \enquote{matOUR.periodic} in periodic form, for restarting or other
  manipulations.
\end{itemize}

It also produces a file tdph.log containing these columns:

\begin{enumerate}
\def\labelenumi{\arabic{enumi}.}
\tightlist
\item
  The iteration number
\item
  The chi square
\item
  and onward: the phonon parameters.
\end{enumerate}

If using minimization = \enquote{ph+zstar}, the zstar minimization will
be at the end of the file (with a different number of columns)

\hypertarget{common-input-cards}{%
\subsection{Common input cards}\label{common-input-cards}}

\hypertarget{qpoints}{%
\subsubsection{QPOINTS}\label{qpoints}}

The \enquote{QPOINTS} card instruct the code to start reading a list of
nq q-points, nq has been previously entered in the namelist. On the same
line as QPOINTS thre optional keywords can be specified:

\hypertarget{cartesian-default}{%
\paragraph{cartesian (DEFAULT)}\label{cartesian-default}}

The points will be specified on Cartesian axis in units of 2π/alat,
where alat is the lattice parameter. This is the default behavior of the
code and is equivalent to not specifying any keyword.

\hypertarget{crystal}{%
\paragraph{crystal}\label{crystal}}

The points are going to be specified on the basis of the reciprocal
lattice vectors.

\begin{center}\rule{0.5\linewidth}{0.5pt}\end{center}

In the default, \enquote{cartesian} and \enquote{crystal} cases, the
code will now start reading the q-points, one per line. After the
q-point cooordinates, each line can optionally include an integer
number, np, which instructs the code to make a straight path from the
previous point to this one formed by np+1 points.

The code will also compute the total length of the path along the
q-points. This length is printed in the output file (usually, 2nd
column) and is useful for plotting.

You can use the special value -1 for np, to reset the length of the path
at a certain point. The d3\_lw.x code will also perform a special action
depending on the context:

\begin{itemize}
\tightlist
\item
  When doing a linewidth calculation: print an empty line (useful for
  gnuplot 3D plots).
\item
  When doing a spectral function calculation: continue writing to a new
  output file, which will have "\_pN" appended to its name; N is a
  number increasing at each path reset.
\end{itemize}

If two consecutive points in the list are equivalent minus a G-vector,
and if the latter has no np, or np=1, then the two points are added to
the list with the same path length. This allows one to jump between
equivalent point without having a discontinuity in the plot.

Another special value is the special value np=0 skips the point, but
puts in in the list as null \enquote{previous point}. I.e. the path will
continue from the next point, as if the null point was the previous one.
This allows to introduce a discontinuity in the path without resetting
the length.

\hypertarget{grid}{%
\paragraph{grid}\label{grid}}

Only used by the d3\_lw.x code when the calculation type is
\enquote{grid}. On the next line the code will try to read three integer
numbers, NQX, NQY and NQZ, which will define the dimension of the grid.

\hypertarget{bz}{%
\paragraph{bz}\label{bz}}

Equivalent to \enquote{grid} but the grid will be translated in the
Brillouin zone (i.e.~the Wigner-Seitz cell of the reciprocal lattice),
points on the boundaries will be duplicated to all the equivalent ones.
If you specify grid or bz, you can add, on the same line, a vector by
which the grid will be shifted. The vector is in units of half the grid
spacing, i.e.~a value of 1 indicates the standard Monkhorst-Pack shift,
you can use any value.

\hypertarget{xsf-or-bxsf}{%
\paragraph{xsf or bxsf}\label{xsf-or-bxsf}}

Only used when doing final state decomposition. Produces a filesuitable
for 3D plotting with XCrysDen, in the xsf or bxsf format. The first is
more flexible, but the second is more suitable for Brillouin-zone plot.
The code will try to read three integer numbers, NQX, NQY and NQZ, note
that the grid size will actually be NQX+1, for xsf, according to
XCrysDen conventions.

\hypertarget{plane}{%
\paragraph{plane}\label{plane}}

Specify a plane of q-points, it will read four lines:

\begin{verbatim}
n1 n2
e1x e1y e1z
e2x e2y e2z
q0x q0y q0z
\end{verbatim}

Which are the grid size (n1, n2) two vectors defining the surface
(e1,e2) and the origin (q0). The first point in the surface is q0, the
last is q0+e1+e2

\hypertarget{examples}{%
\paragraph{Examples}\label{examples}}

Select the Γ point (000) and the X point (001) in a cubic lattice:

\begin{verbatim}
QPOINTS
 2
0.0 0.0 0.0
0.0 0.0 1.0 
\end{verbatim}

Make a path K→ Γ → M in a hexagonal lattice, the path will include a
total of 41 points:

\begin{verbatim}
QPOINTS
 3
0.0 0.33333 0.33333
0.0 0.0 0.0   20
0.0 0.0 0.5   20
\end{verbatim}

Make a line K→ Γ and then M→ Γ (in a hexagonal lattice), the path will
include a total of 82 points:

\begin{verbatim}
QPOINTS crystal
 4
0.33333 0.33333 0.0
0.0     0.0     0.0   20
0.0     0.5     0.0   0
0.0     0.0     0.0   20
\end{verbatim}

Make a 10 × 10 × 10 points grid

\begin{verbatim}
QPOINTS grid
10 10 10
\end{verbatim}

Make a 10 × 10 × 1 points grid, shifted in the plane 10 × 10, but not
along the third direction.

\begin{verbatim}
QPOINTS grid 1 1 0
10 10 1
\end{verbatim}

Make three parallel lines of 101 points each, reset the path length
after each line

\begin{verbatim}
QPOINTS
 6
0 0   0
0 0   1  100
0 0.1 0  -1
0 0.1 1  00
0 0.2 0  -1
0 0.2 1  100
\end{verbatim}

\hypertarget{configs}{%
\subsubsection{CONFIGS}\label{configs}}

The \enquote{CONFIGS} card instructs the code to start reading a list of
sigma (in cm-1)/temperature (in K) configurations. Sigma can be the
width of a Gaussian smearing or the regularization of the self-energy,
depending on the type of calculation. Configurations can be specified as
a list (default), or as a matrix. In the former case a simple list is
expected; in the latter case the code reads two lists: one of sigma and
one of T and generates all possible couples. list

If you specify \enquote{CONFIGS list} or just \enquote{CONFIGS}, the
number of configurations \enquote{nconf} is expected on the next line
(deprecated: you can also use the nconf variable in the namelist).
Afterwards, the code will read nconf lines, and expects a couple
\enquote{sigma T} on each line. If on a certain line only one number is
present, the code assumes it to be a temperature and it reuses the
previous value of the delta.

\emph{matrix} After \enquote{CONFIGS matrix} the code will scan for the
keyword \enquote{T} (\enquote{sigma}), followed, on the same line, by
the number of temperatures (sigmas), the default value of 1 is taken if
missing. The code will the read, on the following lines, the required
values of temperature (sigma). See also the example below.

\hypertarget{examples-1}{%
\paragraph{Examples}\label{examples-1}}

Select three configurations of increasing values of sigma and
temperature:

\begin{verbatim}
CONFIGS
 3
1.0  100
2.0  200
3.0  300
\end{verbatim}

Select five different temperatures, all with the same smearing of 1cm-1:

\begin{verbatim}
CONFIGS
 5
1.0     0
    100
    200
    300
    400
\end{verbatim}

Build thirty configurations, from 10 values of temperature (K) and 3 of
smearing (cm-1):

\begin{verbatim}
CONFIGS matrix
  T  10
 10.  20.  50. 100.  200.
300. 400. 500. 750. 1000.
 sigma  3
2.
3.
4.
\end{verbatim}

Undocumented: using a negative or zero value for the smearing. This
features can change at any time in the future, possibly to become a
proper input variables. A negative value of the smearing can activate a
few \enquote{hidden} features of the d3\_lw.x code. In particular, when
doing a \enquote{lw full} or \enquote{final state} calculation, setting
the smearing to zero will use the static formula for the linewidth,
where the energy denominators are just ω1+ω2 and ω1-ω2, the latter
becomes a derivative (of the Bose-Einstein distribution) when the modes
2 and 3 are degenerate. A negative value of sigma, will also use the
static limit but with a regularization of the denominator. The static
limit is not currently implemented for the spectral function, and it
does not make sense for the \enquote{lw imag} calculation.

\hypertarget{isotopes}{%
\subsubsection{ISOTOPES}\label{isotopes}}

After this card the code will read the information about the isotopic
composition of all elements present in the calculation. This list will
only be read if \enquote{isotopic\_disorder} is set to true in the
namelist. Note that specifying the isotopes is not necessary, if this
card is omitted the natural isotopic concentration will be used.

Each element, identified by its name, must appear in the same order as
in the file of the force constants. The isotopic composition can be
specified in several different ways.

An isotopically pure element can be specified either by atomic number of
the isotope or by setting manually its mass in atomic mass units
(i.e.~mass of 12C=12):

\begin{verbatim}
Xx N natm
\end{verbatim}

Where Xx is the name of the element (e.g., H, Na, Cl), \enquote{N} is a
keyword and natm is the atomic number

\begin{verbatim}
Xx M matm
\end{verbatim}

Where Xx is the name of the element (e.g., H, Na, Cl), \enquote{M} is a
keyword and matm is the atomic mass

If you wish to use the natural isotope concentration for an element you
can use the following syntax:

\begin{verbatim}
Xx natural
\end{verbatim}

Where Xx is the name of the element and \enquote{natural} is a keyword.
The natural isotope concentration is stored in the file
nist\_isotopes\_db.f90 of the thermal2 distribution. It is obtained from
the NIST online database, available on the NIST website7.

If you wish to set the isotope concentration by hand you can use the
following syntax

\begin{verbatim}
Xx isotopes niso
 m1 c1
 …
 mniso cniso
\end{verbatim}

Where Xx is the name of the element, \enquote{isotopes} is a keyword and
niso is the number of isotopes. You will then list the mass and
concentration of each isotopes, on niso separate lines.

Finally, you can also specify directly the average mass for the element
and its variance, which are the only two quantities that are actually
used in the linewidth and thermal conductivity calculation. You can use
the following syntax:

\begin{verbatim}
Xx manual gm gs
\end{verbatim}

Where Xx is the element name, \enquote{manual} is a keyword, gm is the
average mass and gs is its variance (both in Dalton units).

\hypertarget{examples-2}{%
\paragraph{Examples}\label{examples-2}}

In the following example we will set Hydrogen isotopic concentration to
1) pure one-proton H 2) pure Deuterium 3) 50\% H and 50\% D 4) Hydrogen
natural concentration 5) set gm and gs manually

\begin{verbatim}
H N 1
H M 2.01410178
H isotopes 2
    1.00000000 0.5
    2.01410178 0.5
H natural
H manual 1.02  0.01
\end{verbatim}

\hypertarget{input-examples}{%
\section{Input Examples}\label{input-examples}}

In the next sections you will find some example input files for the main
codes of thermal2. You will find more examples, included a

\hypertarget{d3_lw.x-examples}{%
\subsection{d3\_lw.x examples}\label{d3_lw.x-examples}}

\hypertarget{computing-the-linewidth-of-a-single-point}{%
\subsubsection{Computing the linewidth of a single
point}\label{computing-the-linewidth-of-a-single-point}}

This example input would compute the linewidth of q-point (1/3 1/3 0) in
crystal (aka fractional) coordinates for 15 different values of
temperature.

\begin{verbatim}
&lwinput 
 calculation = 'lw real' 
 prefix="lw_rnozp" 

 ! file_mat2 = '../1l.FILDYN/mat2R.120ry.4.nozeu' 
 file_mat2 = '../FILDYN/mat2R.8_nozeu' 
 file_mat3 = '../FILD3DYN/mat3R.xxx_asr15_sparse' 
 outdir    = './' 
 asr2 = 'simple' 
 nconf = 15
 nk = 200,200,1 
 nq = 1
 sort_freq = "overlap"
/ 
CONFIGS 
2.0     0 
       10 
       50 
       100 
       150 
       200 
       250 
       300 
       350 
       400 
       500 
       600 
       700 
       800 
       900 
QPOINTS crystal 
0.33333333333 0.333333333333 0.0
\end{verbatim}

\hypertarget{testing-convergence-with-smearing-and-grid}{%
\subsubsection{Testing convergence with smearing and
grid}\label{testing-convergence-with-smearing-and-grid}}

The following example is useful for testing the convergence of the
linewidth calculation as a function of the smearing value. The input
should be run several times replacing NK with increasing values. You can
then collect the linewidth of a band and plot a curve for every smearing
as a function of the grid size.

\begin{verbatim}
&lwinput 
 calculation = 'lw imag' 
 prefix = 'lw_NK' 
 file_mat2 = '../FILDYN/mat2R.100Ry.4' 
 file_mat3 = '../FILD3DYN/mat3R_asr_sparse' 
 outdir    = './' 
 asr2 = 'simple'  
 nconf = 7
 nk = NK, NK, NK
 nq = 1
/ 
CONFIGS 
 0.1 300
 0.2 300
 0.5 300
 1.0 300
 2.0 300
 5.0 300
10.0 300
QPOINTS 
0 0 0
\end{verbatim}

You can then extract the data with a script like this:

\begin{verbatim}
#!/bin/bash
for smr in 0.1 0.2 0.5 1.0 2.0 5.0 10.0;do
    grep "^ *1 " lw_*s${smr}.out|\
    awk '{print $1,$(NF-2),$(NF-1),$NF}'|\
    sed -re 's/lw_|_T|_s|.out|:/  /g' |\
    sort -k 1n > s${smr}.dat;
done 
\end{verbatim}

And proceed to plot the curves included in the files s0.1.dat \ldots{}
s10.0.dat, e.g.~with gnuplot.
\includegraphics[width=0.7\textwidth,height=\textheight]{images/conv.png}

You see that larger values of smearing converge with smaller grids but
not to exactly the same value as smaller smearing. Using a smearing of
the same order of magnitude as the average linewidth of the acoustic
band is a good idea as the smearing can be identified with the intrinsic
uncertainty of the phonon energy. An ideal procedure would use the
actual linewidth for the smearing, but this should be done iteratively
which would be much more computationally expensive. This scheme may
eventually be implemented in the code if there is enough request.

In brief: do not take a value too small for the smearing, something of
the order of a cm-1 is usually fine.

\hypertarget{compute-the-linewidth-along-a-path-in-the-bz}{%
\subsubsection{Compute the linewidth along a path in the
BZ}\label{compute-the-linewidth-along-a-path-in-the-bz}}

\begin{verbatim}
&lwinput 
 calculation = 'lw full'
 file_mat2 = 'mat2R' 
 file_mat3 = 'mat3R.xxx_asr_sparse' 
 outdir    = './' 
 asr2 = 'simple' 
 nconf = 15 
 nk = 100,100,10 
 nq = 5 
 sort_shifted_freq = .true. 
/ 
CONFIGS 
5.0     0 
       10 
       50 
       100 
       150 
       200 
       250 
       300 
       350 
       400 
       500 
       600 
       700 
       800 
       900 
QPOINTS 
0 0 0 
0.5 0.0 0.0 100 
0.33333333333 0.333333333333 0.0 74 
0.0 0.0 0.0 94 
0.0 0.0 0.5 25
\end{verbatim}

\hypertarget{compute-the-spectral-function-along-a-path}{%
\subsubsection{Compute the spectral function along a
path}\label{compute-the-spectral-function-along-a-path}}

Spectral function calculations are not as well optimized as linewidth
ones and can be much more memory intensive. If you are running out of
memory, do not use more configuration that necessary, especially if
running on massive parallel machines such as IBM BlueGENE.

\begin{verbatim}
&lwinput 
 calculation = 'spf'
 file_mat2 = 'mat2R' 
 file_mat3 = 'mat3R.xxx_asr_sparse' 
 outdir    = './' 
 asr2 = 'simple' 
 nconf = 1
 nk = 100,100,10 
 nq = 5 
 ne = 500
 de = 1.0
/ 
CONFIGS 
5.0     300 
QPOINTS 
0 0 0 
0.5 0.0 0.0 100 
0.33333333333 0.333333333333 0.0 74 
0.0 0.0 0.0 94 
0.0 0.0 0.5 25
\end{verbatim}

You can obtain a color plot like those of
ref.~\protect\hyperlink{ref5}{5} with this gnuplot script:

\begin{verbatim}
set log cb 
splot 'spf_full_T300_s0.3.out' u 1:2:3 w image
\end{verbatim}

\hypertarget{compute-the-final-state-decomposition}{%
\subsubsection{Compute the final state
decomposition}\label{compute-the-final-state-decomposition}}

Final state can be decomposed over the energy, giving a kind of DOS of
the final state, over the q-points or both. The next example decomposes
only over energy.

\begin{verbatim}
&lwinput 
 calculation = 'final' 
 file_mat2 = 'T_295K/mat2R' 
 file_mat3 = '../mat3R_asr' 
 outdir    = 'T_295K/' 
 asr2 = "simple"
 nconf = 3 
 nk =  20, 20, 20 
 nq = 1 ! actually unused
 e_initial = 490.80 
 q_initial = -0.5, 0.0, 0.0 
 ne =201 
 de = 3.5 
 e0 = 0. 
/ 
CONFIGS 
 5      295 
10      295 
20      295 
QPOINTS 
0.0 0.0 0.0
\end{verbatim}

Final state over a high-symmetry path in the brillouin zone.

\begin{verbatim}
&lwinput
  calculation = 'final'
  prefix="final"
  file_mat2 = 'mat2R'
  file_mat3 = 'mat3R.asr.sparse'
  outdir    = './LW/'  
  asr2 = 'simple'   
  nk =  31, 31, 31
  xk0 = 1,1,1
  sort_freq = 'overlap'
  ne = 2000
  de = 0.5

  e_initial = 682.  ! a
  q_initial = 0.0000010000,    0.0000000000,    0.0000000000 
  q_resolved = .true.
  sigmaq= 0.02
/
 CONFIGS
  12
 5.0 300
     400
     500
     600
     700
     800
     900
     1000
     1100
     1200
     1300
     1400
 QPOINTS 
5
0.0000000000    0.0000000000    0.0000000000 1 Γ
1.0000000000    0.0000000000    0.0000000000 20 X
1.0000000000    1.0000000000    0.0000000000 20 X
0.0000000000    0.0000000000    0.0000000000 50 Γ
0.5000000000    0.5000000000    0.5000000000 20 L
\end{verbatim}

\begin{figure}
\centering
\includegraphics[width=0.7\textwidth,height=\textheight]{images/final.png}
\caption{final state over a path in the BZ}
\end{figure}

Final state as a 3D density map, plotted with XCrysDen. Note that the
40×40×40 grid used for plotting does not have to be the same as the nk
grid used to compute the linewidth. Best result are obtained when it is
considerably coarser, as here. Using a finer grid for plotting than for
computing will just result in \enquote{balls} appearing at the
integration points.

\begin{verbatim}
&lwinput
  calculation = 'final'
  prefix="G394-50K-111"
  file_mat2 = 'FILDYN/mat2R'
  file_mat3 = 'FILD3DYN/mat3R.asr.sparse'
  outdir    = './FS/'
  asr2 = 'simple'
  nk = 200,200,200
  grid_type = 'random'
  q_initial = 0.05, 0.00, 0.00
  nu_initial = 3

  ne = 650
  de = 0.5

  q_resolved = .false.
  q_summed = .true.
  sigmaq= 0.05

/
 CONFIGS 
 1
 2.0  50
QPOINTS bxsf
 40 40 40
\end{verbatim}

\begin{figure}
\centering
\includegraphics[width=0.7\textwidth,height=\textheight]{images/LA-0.05-111.png}
\caption{final state in 3D with XCrysDen}
\end{figure}

\hypertarget{color-map-plot-of-the-lw-in-the-bz-2d-systems}{%
\subsubsection{Color-map plot of the LW in the BZ (2D
systems)}\label{color-map-plot-of-the-lw-in-the-bz-2d-systems}}

In bidimensional materials it can be interesting to plot the linewidth
of a certain band over the entire Brillouin zone. In a 3D materials, it
is more tricky to get a nice picture, you can use the
\protect\hyperlink{qpoints}{option \enquote{plane}} to plot a section of
the unit cell.

\begin{verbatim}
&lwinput 
 calculation = 'lw full' 
 prefix="lw_2d" 
 file_mat2 = '../../1l.FILDYN/mat2R.4.vac.nozeu' 
 file_mat3 = '../../1l.FILD3DYN/mat3R.xxx_asr_sparse' 
 outdir    = './' 
 asr2 = 'simple'
 nconf = 1 
 nk = 20,20,1 
 sort_shifted_freq = .false. 
/ 
CONFIGS 
1.0    300
QPOINTS bz 
 50 50 1
\end{verbatim}

You can then plot the resulting file with a gnuplot command like this
(you will have to tune the pointsize and of course select the correct
band to plot):

\begin{verbatim}
p 'lw_2d.50x50x1@bz_T300_s10.out' u 3:4:12 w p palette pt 7 pointsize 1.7 not
\end{verbatim}

This is an example from graphene (done using 128x128 BZ-centered
points):

\begin{figure}
\centering
\includegraphics[width=0.7\textwidth,height=\textheight]{images/graphene-bz.png}
\caption{graphene bz}
\end{figure}

\hypertarget{d3_tk.x-examples}{%
\subsection{d3\_tk.x examples}\label{d3_tk.x-examples}}

\hypertarget{compute-the-sma-solution-of-the-bte}{%
\subsubsection{Compute the SMA solution of the
BTE}\label{compute-the-sma-solution-of-the-bte}}

The following input file computes the thermal conductivity in the
Single-Mode Approximation. This example uses a finer grid of 24x24x24
q-points to compute the phonon lifetime on a grid of 12x12x12 points,
which is integrated to compute the thermal conductivity. Furthermore,
the natural isotopic distribution of Silicon is used.

\begin{verbatim}
&tkinput 
 calculation = 'sma' 
 prefix="tk"
 file_mat2 = './mat2R' 
 file_mat3 = './FILD3DYN/mat3R.asr.sparse' 
 outdir    = './' 
 asr2 = 'simple' 
 nconf = 1
 nk    = 12,12,12
 nk_in = 24,24,24
 casimir_scattering=.false. 
 isotopic_disorder =.true.
/ 
CONFIGS 
5.0    10 
ISOTOPES 
 Si natural
\end{verbatim}

\hypertarget{compute-the-exact-solution-of-the-bte}{%
\subsubsection{Compute the exact solution of the
BTE}\label{compute-the-exact-solution-of-the-bte}}

The following input file computes the thermal conductivity solving the
BTE by functional minimization. This example also includes border
scattering on grains of 1 micrometer average size.

\begin{verbatim}
&tkinput 
 calculation = 'cgp' 
 prefix="tk"
 file_mat2 = './mat2R' 
 file_mat3 = './FILD3DYN/mat3R.asr.sparse' 
 outdir    = './' 
 asr2 = 'simple' 
 nconf = 9
 nk = 24,24,24
 niter_max = 100
 thr_tk = 1.d-4
 casimir_scattering=.true. 
 sample_length_mu = 1.0 
/ 
CONFIGS 
5.0    10 
       20 
       50 
       100 
       200 
       300 
       500 
       1000 
       2000
\end{verbatim}

\hypertarget{bibliography}{%
\section{Bibliography}\label{bibliography}}

A list of essential bibliographic references follows. When using this
code you are not required to cite all of these work, however the authors
would greatly appreciate if you cited references 1-4 where the theory
and implementation of this work is described in detail.

In particular ref. 1 describes the calculation of 3rd order dynamical
matrices using DFPT, and the calculation of ph-ph linewidth; ref. 2
describes the iterative solution of the BTE. In ref. 9 the formula used
to compute the phonon-phonon scattering is derived. See also refs. 10-14
for previous methods, and for the inclusion of isotopic and border
scattering. Ref. 3 and 4 include important applications and analysis of
the BTE exact solution. Reference 5 describes the calculation of phonon
spectral functions and the combination of 3rd and 4th order. References
6-8 contain the foundation of the theory used. The geometric terms used
to include finite-size effects are described in detail in ref. 15.

\begin{enumerate}
\def\labelenumi{\arabic{enumi}.}
\tightlist
\item
  L. Paulatto, F. Mauri, and M. Lazzeri, Phys. Rev.~B 87, 214303 (2013).
\item
  G. Fugallo, M. Lazzeri, L. Paulatto, and F. Mauri, Phys. Rev.~B 88,
  045430 (2013).
\item
  A. Cepellotti, G. Fugallo, L. Paulatto, M. Lazzeri, F. Mauri, N.
  Marzari, Nature communications 6, 6400 (2015)
\item
  G. Fugallo, A. Cepellotti, L. Paulatto, M. Lazzeri, N. Marzari, F.
  Mauri, Nano letters 14 (11), 6109-6114 (2014)
\item
  Lorenzo Paulatto, Ion Errea, Matteo Calandra, and Francesco Mauri,
  Phys. Rev.~B 91, 054304 (2015)
\item
  M. Calandra, M. Lazzeri, F. Mauri, Physica C 456, 38--44 (2007)
\item
  Michele Lazzeri and Stefano de Gironcoli, Phys. Rev.~Lett. 81, 2096
  (1998)
\item
  N. Bonini, M. Lazzeri, N. Marzari, and F. Mauri, Phys. Rev.~Lett. 99,
  176802 (2007)
\item
  J. Garg, N. Bonini, B. Kozinsky et al., Phys. Rev.~Lett. 106, 045901
  (2011)
\item
  Joseph Callaway, Phys. Rev.~113, 1046
\item
  M. Omini, A. Sparavigna, Il Nuovo Cimento, 19, 1537 (1997)
\item
  M. Omini and A. Sparavigna, Phys. Rev.~B 53, 9064 (1996)
\item
  M. Omini, A. Sparavigna, Physica B: Condensed Matter 212, 101-112,
  (1995)
\item
  M. Omini and A. Sparavigna, Phys. Rev.~B 61, 6677 (2000)
\item
  L. Paulatto, D Fournier, M Marangolo, M Eddrief, P Atkinson, M
  Calandra. Phys. Rev.~B 101 (20), 205419 (2020)
\item
  R. Bianco, I. Errea, L. Paulatto, M. Calandra, and F. Mauri. Phys.
  Rev.~B 96, 014111 (2017)
\item
  Simoncelli, Marzari, Mauri, Nature Physics volume 15, 809--813 (2019)
\item
  Simoncelli, Marzari, Mauri, Phys. Rev.~X 12, 041011 (2022)
\end{enumerate}

\hypertarget{change-log}{%
\section{Change Log}\label{change-log}}

See the changelog file in the \enquote{Doc} subdirectory.

\end{document}
